% 从旧数学附录迁出的可持久化线段树说明；不新增算法。
\section{可持久化区间拼接的懒标记}\label{knowledge-persistent-concat}
代码见第~\pageref{compact-PersistentRange}~页。从全零版本开始，同一对象内版本共享坐标域；
标记永久化：$\mathrm{sum}=\mathrm{left.sum}+\mathrm{right.sum}+\mathrm{lazy}\cdot\mathrm{len}$，查询只读。
拼接把目标版本某半开区间替换为源版本同一区间；递归参数 offset 下降时更新为
$\mathrm{offset}+\mathrm{source.lazy}-\mathrm{destination.lazy}$，完整覆盖时复制源节点并把 offset 加到 lazy 与 sum（offset 为零直接共享）。
修改、拼接 $O(\log n)$ 新节点，查询 $O(\log n)$ 零新节点。不支持错位平移粘贴，不回收节点。
和、差、标记与长度乘积须在 long long 内；节点数与版本数小于 \texttt{INT\_MAX}。
题目时间戳可回退时，用 ver[t] 映射时间戳到模板版本，不能把 add 返回的追加编号当时间戳。

\section{杭州 2023 K：跨版本拼接应用}\label{knowledge-hangzhou-k}
0-based，$F_i(r)$ 为 $[i,r]$ 剩余牌数，$p$ 为 $i$ 之后第一张与 $a_i$ 相同的牌（无则 $p=n$）：
\[
 F_i(r)=\begin{cases}
 1+F_{i+1}(r),&r<p,\\
 F_{p+1}(r),&r\ge p.
 \end{cases}
\]
从 $i=n-1$ 向前：复制版本 $i+1$，$[i,p)$ 加 1，$[p,n)$ 拼接自版本 $p+1$。
询问 $[l,r]$ 读版本 $l-1$ 的坐标 $r-1$；强制在线。构建 $O(n\log n)$，查询 $O(q\log n)$。
